% 第 6 章 IC Validator Explorer DRC
\bichapter{IC Validator Explorer DRC}{IC Validator Explorer DRC}
\label{chap:expdrc}

本章说明 IC Validator Explorer DRC 的功能。

IC Validator Explorer DRC 包括以下各节：
\begin{itemize}
  \item IC Validator Explorer DRC 简介
  \item Explorer DRC 模式类型
  \item 配置 Explorer DRC
  \item Explorer DRC 输出文件
  \item 使用 Explorer DRC 调试
  \item Explorer DRC 许可证
\end{itemize}

% ============================================================
\section{IC Validator Explorer DRC 简介}
\subsection*{Introduction to IC Validator Explorer DRC}

在设计早期调试 DRC 违例可能既复杂又耗时。例如，复杂的 DRC 规则（多重图形、delta voltage 等）运行时间很长，而在仍存在其他更简单违例时，这些检查往往并无实际收益。
可用一套基本设计检查来调试典型的早期设计问题：查找 cell 重叠、不当 fill、布线问题等并不需要复杂规则。IC Validator Explorer DRC 会从晶圆厂 runset 中自动确定这些基本规则，并额外加入方法学规则与调试能力，从而显著缩短周转时间（turnaround time）。
此外，当某些设计层触发过多违例时，IC Validator Explorer DRC 会检测并取消选择相关的慢速规则。这种对检查的动态选择与执行，可比传统的完整 DRC 运行获得更快的周转时间。

% ============================================================
\section{Explorer DRC 模式类型}
\subsection*{Types of Explorer DRC Modes}

Explorer DRC 有三种模式：
\begin{itemize}
  \item Explorer DRC Standalone 模式\\
  该模式运行从晶圆厂 runset 中自动选出的设计规则，以及 IC Validator 额外加入的方法学检查。\\
  适用于存在违例、需要快速检查且计算资源有限的设计。

  \item Explorer DRC Tiers 模式\\
  Explorer DRC tiers 提供系统化框架，优先执行并报告高优先级检查。在 tiers 模式下，检查会完整执行，不以性能为主要取舍：所有规则按 tier 顺序选择并执行，直至指定的 tier。\\
  使用该模式时，性能并非主要考虑因素。\\
  部分 tier 保留给 Explorer 优先级规则与方法学检查；多数 tier 可由用户定制。更多信息见表~\ref{tab:explorer-tiers}。\\
  Tiers 按所赋优先级依次执行。编号较小的 tier 优先级最高。位于 Tier~0 与 Tier~10 之间、尚未定义归属的检查，在 Tier~10 之后作为「All remaining rules」tier 执行。每个 tier 完成后，结果会在运行过程中更新。\\
  \textbf{Tiers 模式运行示例}\\
  在全部 Tiers 上运行 Explorer DRC auto 模式，并指定用户定义的 Tiers，使用 \cmd{-explorer\_tiers\_file mytier.txt}：
\begin{lstlisting}
% icv -explorer auto -explorer_tiers_file mytier.txt runsetfile
\end{lstlisting}

  \item Explorer DRC Auto 模式\\
  与 tiers 模式类似，该模式同样使用 Explorer DRC tiers，但还具备停止慢速规则的能力。\\
  在 auto 模式下，若发现存在违例的设计层，相关规则若其函数运行缓慢则会被停止。这样可利用较早 tier 的结果，阻止后续函数/规则运行过久（tier~0 与 tier~1 的规则绝不会被提前停止）。\\
  该模式的性能取决于 DRC 错误数量。DRC 错误越多，auto 模式越可能取消选择规则并提前退出，从而加快周转。\\
  当设计中的违例数量未知、并希望尽可能完成更多检查时，使用 auto 模式。若设计干净，auto 模式不会停止任何检查，并产生完整签核结果。
\end{itemize}

\subsection{执行 Explorer DRC}
\subsubsection*{Executing Explorer DRC}

要执行 Explorer DRC，在命令行设置：
\begin{lstlisting}
-explorer standalone|auto[:TierNumber]|tiers[:TierNumber]
\end{lstlisting}
其中 \cmd{TierNumber} 指定要执行到的最高 tier。将执行从最低到（含）\cmd{TierNumber} 的全部 tiers。若未指定 \cmd{TierNumber}，则所有检查均可作为执行候选。
有关配置 IC Validator Explorer DRC 的更多信息，见下一节。

\subsubsection{Explorer DRC 模式运行示例}
\paragraph*{Examples of Running in Explorer DRC Mode}

运行 Explorer DRC standalone 模式：
\begin{lstlisting}
% icv -explorer standalone runsetfile
\end{lstlisting}

运行 Explorer DRC auto 模式直至 Tier~4：
\begin{lstlisting}
% icv -explorer auto:4 runsetfile
\end{lstlisting}

表~\ref{tab:explorer-tiers} 给出不同 Explorer DRC tiers、预含检查以及可定制类别。

\begin{longtable}{@{}p{0.08\textwidth} p{0.16\textwidth} p{0.14\textwidth} p{0.32\textwidth} p{0.20\textwidth}@{}}
\caption{Explorer Tiers}\label{tab:explorer-tiers}\\
\toprule
\textbf{Tier} & \textbf{用户可否向该 tier 添加检查？} & \textbf{用户可否命名该 tier？} & \textbf{预含检查} & \textbf{可选检查} \\
\midrule
\endfirsthead
\toprule
\textbf{Tier} & \textbf{用户可否向该 tier 添加检查？} & \textbf{用户可否命名该 tier？} & \textbf{预含检查} & \textbf{可选检查} \\
\midrule
\endhead
\bottomrule
\endfoot
0 & 否 & 否 & 方法学警告、fill overlap 诊断、优先级规则 & N/A \\
1 & 是 & 否 & 无 & 优先级规则（用户指定） \\
2 & 是 & 是 & 基于电压的诊断（集合~1） & 用户指定规则 \\
3 & 是 & 是 & 基于电压的诊断（集合~2） & 用户指定规则 \\
4--10 & 是 & 是 & 无 & 用户指定规则 \\
N/A & N/A & N/A & 全部剩余规则（All remaining rules） & N/A \\
\end{longtable}

% ============================================================
\section{配置 Explorer DRC}
\subsection*{Configuring Explorer DRC}

Explorer DRC 模式的命令行选项见表~\ref{tab:explorer-drc-cli}。

\begin{longtable}{@{}>{\raggedright\arraybackslash\ttfamily}p{0.36\textwidth} p{0.56\textwidth}@{}}
\caption{Explorer DRC 模式命令行选项}\label{tab:explorer-drc-cli}\\
\toprule
\textbf{\textrm{选项}} & \textbf{\textrm{说明}} \\
\midrule
\endfirsthead
\toprule
\textbf{\textrm{选项}} & \textbf{\textrm{说明}} \\
\midrule
\endhead
\bottomrule
\endfoot
-explorer standalone\allowbreak|auto[:TierNumber]\allowbreak|tiers[:TierNumber]
  & 以 standalone、auto 或 tiers 模式运行 Explorer。\cmd{TierNumber} 对 auto 与 tiers 模式可选，用于将运行限制到指定的 \cmd{TierNumber}；未指定时，所有 tiers 均为执行候选。 \\
-explorer\_tiers\_file path\_to\_file
  & 指定用于定制用户定义 tiers 并传入 Explorer DRC 运行的文件。该文件中可使用 \cmd{ADD\_RULES\_TO\_TIER} 与 \cmd{NAME\_TIER} 指令。 \\
-explorer\_script path\_to\_file
  & 指定在 Tier~0 与 Tier~1 完成后用于追加资源的文件。可先以较少资源启动 auto 运行，并在 Tier~0 与 Tier~1 完成后动态追加资源，以加快其余检查的周转。 \\
\end{longtable}

命令行选项 \cmd{-explorer\_tiers\_file} 指向可选的脚本文件，用于定制 tiers。该文件示例如下：
\begin{lstlisting}[style=pxl]
ADD_RULES_TO_TIER 1 "*.1:" "*G1*"
ADD_RULES_TO_TIER 2 "*.S.*" "**DN*"
NAME_TIER 4 "Handoff Rules"
ADD_RULES_TO_TIER 4 "*delta*"
ADD_RULES_TO_TIER 4 "*LU*"
DISABLE_EXPLORER.OVERLAP
\end{lstlisting}

创建该文件时可使用下列指令：
\begin{itemize}
  \item \cmd{ADD\_RULES\_TO\_TIER TierNumber Rule comment \ldots{} Rule comment}\\
  用于允许用户指定规则的 tiers。可按违例注释（violation comments）将规则加入各 tiers，支持字符串匹配。\\
  规则按优先级升序加入各 tiers。若某规则已加入更高优先级的 tier，则不会再加入任何更低优先级的 tier；该规则与 Explorer tiers 文件中的书写顺序无关。\\
  允许跳过或留空的 tiers。当 Explorer DRC 遇到空 tier 时，执行会移至下一优先级 tier。报告按 tier 优先级顺序给出。

  \item \cmd{NAME\_TIER TierNumber TierName}\\
  用于允许用户指定名称的 tiers。该名称用于向 \cmd{LAYOUT\_ERRORS} 文件、\cmd{RESULTS} 文件或 VUE 报告。\\
  允许未命名的 tiers；报告时使用默认 Explorer tier 名称 \cmd{TierNumber}。

  \item \cmd{DISABLE\_EXPLORER.DiagnosticName}\\
  用于禁用 Explorer 专有诊断。支持的取值包括：
  \begin{itemize}
    \item \cmd{EXPLORER.OVERLAP}：禁用默认包含在 Tier~0 中的 fill overlap 诊断。
    \item \cmd{EXPLORER.ASSIGN}：禁用默认包含在 Tier~0 中的基于 assign 层的优先级规则。
    \item \cmd{EXPLORER.VOLTAGE}：禁用默认包含在 Tier~3 与 Tier~4 中的基于电压的诊断。
  \end{itemize}
\end{itemize}

命令行选项 \cmd{-explorer\_script path\_to\_file} 在 Tier~0 与 Tier~1 检查均完成后运行。可以少于完成整次运行所需的资源启动 DRC auto 模式；若后续 tiers 需要更多资源，该选项可在基本检查之后追加资源。脚本文件示例如下：
\begin{lstlisting}
#!/bin/csh -f

echo "Explorer Tier 0 and Tier 1 checks are done. Adding resources and
running remaining Tiers"

# Following command adds host "A" to ongoing run
icv -rd run_details -host_add A
\end{lstlisting}

% ============================================================
\section{Explorer DRC 输出文件}
\subsection*{Explorer DRC Output Files}

Explorer DRC 会创建与 IC Validator DRC 运行类似的 ASCII 文件。有关 DRC 运行文件的更多信息，见「输出文件」一章。
在 \cmd{LAYOUT\_ERRORS} 与 \cmd{RESULTS} 文件中，Explorer 运行的错误按表~\ref{tab:explorer-err-cat} 分类。

\begin{longtable}{@{}>{\raggedright\arraybackslash\ttfamily}p{0.28\textwidth} p{0.64\textwidth}@{}}
\caption{Explorer 错误分类}\label{tab:explorer-err-cat}\\
\toprule
\textbf{\textrm{类别}} & \textbf{\textrm{说明}} \\
\midrule
\endfirsthead
\toprule
\textbf{\textrm{类别}} & \textbf{\textrm{说明}} \\
\midrule
\endhead
\bottomrule
\endfoot
\textrm{Methodology warnings}
  & 指定旨在标记潜在设计或性能问题的 Explorer 方法学警告。 \\
\textrm{Fill overlap diagnostic}
  & 指定被诊断为 fill overlap 问题的 DRC 错误（runset 选定的 fill overlap 规则）。 \\
\textrm{Priority rules}
  & 指定来自 Explorer 优先级检查的 DRC 错误（runset 选定的基本规则）。该类别中某些规则可能带有 \cmd{EXPLORER} 前缀，表示该规则已相对原始 runset 被修改（更快版本）。原始 runset 版本会作为「All remaining rules」类别中的 DRC Errors 运行。 \\
\textrm{Priority Rules (user-defined)}
  & 指定来自用户定义优先级规则的 DRC 错误。 \\
\textrm{Voltage dependent diagnostics}
  & 指定标记多边形或 net 对象上基于 text 的电压冲突的 Explorer 方法学警告；亦包含游离电压 text 与电压检查。 \\
\textrm{Explorer tier TierNumber}
  & 指定来自未命名 tiers 的 DRC 错误。默认将这些错误归类在 Explorer tier、\cmd{TierNumber} 下。可用文件中的 \cmd{NAME\_TIER} 指令修改 tier 名称。 \\
\textrm{All remaining rules}
  & 指定来自归类为预定义或用户定义 tiers 之检查的可用 DRC 错误。 \\
\end{longtable}

除表~\ref{tab:explorer-err-cat} 中的错误分类外，DRC 错误还会按用户定义的类别名称分类（通过文件中的 \cmd{NAME\_TIER} 指令定义）。

\begin{noteBox}
错误类别取决于 Explorer 模式类型，以及该模式下所标记的诊断。
\end{noteBox}

\subsection{Explorer DRC RESULTS 文件}
\subsubsection*{Explorer DRC Results File}

Explorer DRC 创建 \cmd{RESULTS} 文件 \cmd{cell.RESULTS}。该文件提供运行结果，以及本次运行所用的 IC Validator 命令行选项。
\cmd{RESULTS} 文件是分析 Explorer DRC 运行结果的起点。Results summary 部分按表~\ref{tab:explorer-err-cat} 的分类报告 DRC 错误统计，包括不同 Explorer 类别中的规则总数与违例总数。

\subsection{Explorer DRC Errors 文件}
\subsubsection*{Explorer DRC Errors File}

Explorer DRC 创建 \cmd{LAYOUT\_ERRORS} 文件，提供运行期间标记的版图错误；这些错误按表~\ref{tab:explorer-err-cat} 分类。
为使错误报告简洁，\cmd{LAYOUT\_ERRORS} 文件中会抑制错误坐标信息，等效于运行 \cmd{error\_options(report\_error\_details = false)}。信息仍保存在错误数据库（PYDB）中，但仅将违例摘要写入 \cmd{LAYOUT\_ERRORS} 文件。
每个 tier/类别完成后即生成 \cmd{LAYOUT\_ERRORS} 文件，以便在 Explorer DRC 运行过程中查阅。\cmd{icv.log} 会在 tier/类别完成时指示，例如：
\begin{lstlisting}
ICV_Engine run is 5% complete. Elapsed Time=0:07:24
All Priority Rules are finished! 19103703 violations were found.
\end{lstlisting}

\begin{noteBox}
在将错误标记存入错误数据库（PYDB）时，Tier~0 与 Tier~1（Explorer 优先级检查与用户定义优先级检查）会强制施加 \cmd{error\_sampling\_per\_check} 参数所设限制（默认 100 万）。其他 tiers 的限制则基于 runset。有关 \cmd{error\_sampling\_per\_check} 参数的更多信息，见 \textit{IC Validator Reference Manual}。
\end{noteBox}

\subsection{Explorer DRC VUE 文件}
\subsubsection*{Explorer DRC VUE File}

IC Validator Explorer 创建 VUE 文件 \cmd{cell.vue}，用于查看 Explorer DRC 错误详情与 Heat Map。

% ============================================================
\section{使用 Explorer DRC 调试}
\subsection*{Debugging With Explorer DRC}

分析 Explorer DRC 运行时：
\begin{enumerate}
  \item 以 Explorer DRC 运行 IC Validator。
  \item 审查 \cmd{RESULTS} 文件。
  \begin{itemize}
    \item 若结果干净，进入签核 DRC；若不干净，进入步骤~3。
  \end{itemize}
  \item 将 Explorer DRC 的 \cmd{.vue} 文件加载到 VUE，并将版图编辑器附着到 VUE。
  \item 点击 \textbf{Heat Map} 选项卡开始调试。Error Heat Map 是用于可视化错误分布的图形界面。\\
  Heat Maps 对 Explorer DRC 自动启用。更多信息见 \textit{IC Validator VUE User Guide} 中的「Introduction to the IC Validator Error Heat Map」一章。
  \item 与 \cmd{RESULTS}、\cmd{LAYOUT\_ERRORS} 文件类似，Explorer DRC 错误默认按表~\ref{tab:explorer-err-cat} 分类。一般按以下顺序开始调试：方法学警告、诊断、预定义优先级规则、用户定义 tiers，最后到全部剩余规则。\\
  此外，可在 Violation Category 窗口的菜单中选择 Clustered Errors 类别模式。该模式下相同错误会被聚类，从而可一次调试多条规则。有关 clustered errors 的更多信息，见 \textit{IC Validator VUE User Guide} 中的「Introduction to the IC Validator Error Heat Map」一章。
  \item 从各自类别选定规则后，定位错误计数很高的区域（热图中的红色区域），并放大这些密集区域。
  \item 在这些区域中探测热图与版图错误，分析高错误计数的根本原因。探测可包括：
  \begin{itemize}
    \item 放大区域
    \item 高亮这些区域中的错误标记
    \item 选择或取消选择特定规则
  \end{itemize}
  \item 重复步骤~6 与~7，直至穷尽所有类别/聚类。
\end{enumerate}

% ============================================================
\section{Explorer DRC 许可证}
\subsection*{Explorer DRC Licenses}

IC Validator Explorer DRC 需要 Elite 许可。该功能需要下列 IC Validator 许可证：
\begin{itemize}
  \item \cmd{ICValidator-Manager}
  \item \cmd{ICValidator2-GeometryEngine}
\end{itemize}

\begin{seeAlsoBox}
更多信息见本手册「许可证与资源需求」一章。
\end{seeAlsoBox}
